\documentclass[10pt]{article}
\usepackage[a4paper,landscape,margin=1.5cm]{geometry}
\usepackage{fontspec}
\setmainfont{DejaVu Serif}
\setmonofont{DejaVu Sans Mono}[Scale=0.9]
\usepackage{amsmath,amssymb,longtable,booktabs,array,xurl,fancyvrb}
\usepackage[hidelinks]{hyperref}
\renewcommand{\arraystretch}{1.25}
\setlength{\tabcolsep}{4pt}
\setlength{\parindent}{0pt}
\setlength{\parskip}{4pt}
\title{Inequality and Social Contests: manuscript skeleton}
\date{Pass 2}
\begin{document}
\maketitle
Source manuscript for the model claims: firmworkers\_model.pdf, 15 March 2026 (Appendix: The closed N-player economy; Linear Approximation for Risk-Premium).

Headline and concluding claims may rest only on model claims, literature claims and measurement-map rows. Model claims are proved in Lean 4 against Mathlib. The full Lean source is in Appendix~A and the build files in Appendix~B.

\section*{Status vocabulary}
\begin{description}
\item[\texttt{LEAN\_PROVED}] Lean theorem compiles against Mathlib without sorry
\item[\texttt{LEAN\_WRITTEN}] Lean theorem written, not yet compiled
\item[\texttt{SYMPY}] checked by the SymPy script, no Lean proof yet
\item[\texttt{OPEN}] stated in the manuscript, no formal object yet
\item[\texttt{UNDERSPECIFIED}] the manuscript leaves a primitive undefined
\item[\texttt{ILL\_POSED}] not well-defined as printed
\item[\texttt{REFUTED}] contradicted by a proved or checked claim
\item[\texttt{VERBATIM}] quote extracted from the primary text supplied, page given, checked against local text
\item[\texttt{PRIMARY\_WEB}] primary full text read through the web tool, quote not verbatim-checked
\item[\texttt{ABSTRACT}] only the abstract or publisher summary was read
\item[\texttt{SECONDARY}] supported only by a secondary source
\end{description}

\section{Headline claims}
No headline claims are admitted in this pass.

\section{Model claims}
{\small
\begin{longtable}{>{\raggedright\arraybackslash}p{0.8cm} >{\raggedright\arraybackslash}p{5.0cm} >{\raggedright\arraybackslash}p{2.9cm} >{\raggedright\arraybackslash}p{3.9cm} >{\raggedright\arraybackslash}p{4.4cm} >{\raggedright\arraybackslash}p{3.6cm} >{\raggedright\arraybackslash}p{2.5cm}}
\toprule ID & Claim & Assumptions & Manuscript anchor & Evidence & Proof location & Status \\ \midrule \endhead
M1 & $x_i = x_0(1-i)^{-1/a}$ for $1 \le i \le N$ defines the income ladder & $a>0$ & Appendix, condition 3: "pareto income list x0 (1 - i)\textasciicircum{}\{-1/a\} for 1 \textless{}= i \textless{}= N" & Base $1-i=0$ at $i=1$ and $1-i<0$ for $i\ge 2$ & SymPy \nolinkurl{m1_printed_index_ill_posed}, Listing C.1, lines 10--12 & \texttt{ILL\_POSED} \\ \midrule
M2 & $x(q)=x_0(1-q)^{-1/a}$ is strictly increasing in $q$ & $x_0>0$, $a>0$, $q_1<q_2<1$ & repair of M1 with a quantile index & --- & Lean \nolinkurl{Isc.pareto_strictMono_quantile}, Listing A.3 (\nolinkurl{Ladder.lean}), lines 7--15 & \texttt{LEAN\_PROVED} \\ \midrule
M3 & $x(q)$ is strictly decreasing in $a$ at every interior quantile & $x_0>0$, $0<a_1<a_2$, $0<q<1$ & repair of M1 with a quantile index & --- & Lean \nolinkurl{Isc.pareto_strictAnti_tail}, Listing A.3 (\nolinkurl{Ladder.lean}), lines 17--23 & \texttt{LEAN\_PROVED} \\ \midrule
M4 & $x(q_2;a_2)-x(q_1;a_2) < x(q_2;a_1)-x(q_1;a_1)$, so every income gap on the ladder shrinks as $a$ rises & $x_0>0$, $0<a_1<a_2$, $0\le q_1<q_2<1$ & repair of M1 with a quantile index & --- & Lean \nolinkurl{Isc.pareto_gap_strictAnti_tail}, Listing A.4 (\nolinkurl{Gap.lean}), lines 23--45; SymPy \nolinkurl{m4_adjacent_gap_decreasing_in_a}, Listing C.1, lines 15--30 & \texttt{LEAN\_PROVED} \\ \midrule
M5 & Higher $a$ raises the stakes across income levels & ladder of M2 & Appendix, condition 3: "A higher a corresponds to higher-stakes through income levels." & Contradicted by M3 and M4. $x(0.9)/x_0 = 4.642, 3.162, 2.154, 1.585$ at $a=1.5,2,3,5$ & --- & \texttt{REFUTED} \\ \midrule
M6 & The printed cost function charges $(m_k-m_i)^b$ for downward moves and nothing for upward moves & $b>0$ & Appendix, condition 2: "C(mi, mk) = (mk - mi)\textasciicircum{}b for mi \textgreater{} mk, 0 for mi \textless{} mk" & Base $m_k-m_i<0$ on the charged branch, complex at $b=1/2$. Contradicts the next sentence of the manuscript, "no cost is paid for moving from a position of high-skill to low". $m_i=m_k$ unassigned & SymPy \nolinkurl{m6_printed_cost_complex_branch}, Listing C.1, lines 38--40 & \texttt{ILL\_POSED} \\ \midrule
M7 & $C(m_i,m_k)=(m_k-m_i)^b$ if $m_i<m_k$ and $0$ otherwise. Downward and lateral moves cost $0$ & $m_k\le m_i$ & repair of M6 & --- & Lean \nolinkurl{Isc.upCost_down_free}, Listing A.3 (\nolinkurl{Ladder.lean}), lines 27--28 & \texttt{LEAN\_PROVED} \\ \midrule
M8 & $C(m_1,m_2)+C(m_2,m_3)\le C(m_1,m_3)$. Two adjacent upward moves cost no more than the direct jump & $b\ge 1$, $m_1<m_2<m_3$ & repair of M6 & --- & Lean \nolinkurl{Isc.upCost_steps_le_jump}, Listing A.3 (\nolinkurl{Ladder.lean}), lines 30--36 & \texttt{LEAN\_PROVED} \\ \midrule
M9 & $C(m_1,m_3)\le C(m_1,m_2)+C(m_2,m_3)$. The direct jump costs no more than two steps & $0\le b\le 1$, $m_1<m_2<m_3$ & repair of M6 & --- & Lean \nolinkurl{Isc.upCost_jump_le_steps}, Listing A.3 (\nolinkurl{Ladder.lean}), lines 38--45 & \texttt{LEAN\_PROVED} \\ \midrule
M10 & A swap of positions $i,k$ leaves total income $\sum_j x_j$ unchanged & $x:\{1..N\}\to\mathbb{R}$ & Appendix, condition 1: "another worker k out of N workers is randomly selected to swap her job with worker i" & --- & Lean \nolinkurl{Isc.swap_preserves_total}, Listing A.2 (\nolinkurl{Lottery.lean}), lines 13--15 & \texttt{LEAN\_PROVED} \\ \midrule
M11 & Expected net lottery gain $G_i=\frac{S-Nx_i}{N-1}-L$, with $S=\sum_k x_k$ & $k$ uniform over the $N-1$ other workers, risk-neutral valuation. Uniformity is added to the manuscript & Appendix, condition 1 & --- & Lean \nolinkurl{Isc.lotteryGain_closed}, Listing A.2 (\nolinkurl{Lottery.lean}), lines 31--33 & \texttt{LEAN\_PROVED} \\ \midrule
M12 & $\sum_i G_i=-NL$. The lottery is zero-sum in income and negative-sum net of entry costs & as M11, $N\ge 2$ & Appendix, condition 1 & --- & Lean \nolinkurl{Isc.lotteryGain_total}, Listing A.2 (\nolinkurl{Lottery.lean}), lines 35--44 & \texttt{LEAN\_PROVED} \\ \midrule
M13 & $G_i>0 \iff x_i<\frac{S-(N-1)L}{N}$ & as M11, $N\ge 2$ & Appendix, condition 1 & --- & Lean \nolinkurl{Isc.lotteryGain_pos_iff}, Listing A.2 (\nolinkurl{Lottery.lean}), lines 46--52 & \texttt{LEAN\_PROVED} \\ \midrule
M14 & The swap partner $k$ is drawn from lottery participants &  & Appendix, condition 1: "Worker k does not pay any additional cost other than the fixed lottery cost L" & The manuscript does not say whether $k$ is drawn from participants or from all $N$ workers. If from all $N$, non-participants are swapped without choosing to play & --- & \texttt{UNDERSPECIFIED} \\ \midrule
M15 & $\mathbb{E}[u(\bar y)+u^{\prime}(\bar y)\sigma\varepsilon+\tfrac12u^{\prime\prime}(\bar y)\sigma^2\varepsilon^2]=u(\bar y)+\tfrac12u^{\prime\prime}(\bar y)\sigma^2$ & $\varepsilon$ with finite support, $\mathbb{E}\varepsilon=0$, $\mathbb{E}\varepsilon^2=1$. General distributions not covered & Appendix, Linear Approximation: "E[u(y)] = u(ybar) + 1/2 u′′(ybar) sigma\_y\textasciicircum{}2" & --- & Lean \nolinkurl{Isc.taylor_expectation}, Listing A.2 (\nolinkurl{Lottery.lean}), lines 54--65 & \texttt{LEAN\_PROVED} \\ \midrule
M16 & $\tfrac12u^{\prime\prime}\sigma^2>0\iff u^{\prime\prime}>0$ & $\sigma\ne 0$ & Appendix, Linear Approximation: "positive when u′′ \textgreater{} 0" & --- & Lean \nolinkurl{Isc.premium_pos_iff}, Listing A.2 (\nolinkurl{Lottery.lean}), lines 67--78 & \texttt{LEAN\_PROVED} \\ \midrule
M17 & The lottery variance equals $p(1-\lambda)^2$ with $p<0$ &  & Appendix, Linear Approximation: "sigma\_y\textasciicircum{}2 = c((1 - lambda\_H) m\_H)\textasciicircum{}2 = p\_H (1 - lambda\_H)\textasciicircum{}2 (p\_H \textless{} 0)" & Left side $\ge 0$, right side $<0$ for $\lambda<1$ & SymPy \nolinkurl{m17_variance_sign_conflict}, Listing C.1, lines 43--47 & \texttt{ILL\_POSED} \\ \midrule
M18 & If $\sigma_y=s(1-\lambda)m$ then $\tfrac12u^{\prime\prime}\sigma_y^2=(\tfrac12u^{\prime\prime}s^2m^2)(1-\lambda)^2$. The coefficient $p=\tfrac12u^{\prime\prime}s^2m^2$ varies with income $m$ &  & repair of M17 & --- & Lean \nolinkurl{Isc.premium_scales_with_income}, Listing A.2 (\nolinkurl{Lottery.lean}), lines 80--82 & \texttt{LEAN\_PROVED} \\ \midrule
M19 & The convex region of $u$ lies below a reference income $r$ &  & Appendix, Linear Approximation: "positive when u′′ \textgreater{} 0 (convex region -\textgreater{} below r)" & $r$ is not defined in the manuscript. No functional form of $u$ is given & --- & \texttt{UNDERSPECIFIED} \\ \midrule
M20 & In long-run equilibrium high ranks choose the downgrade and low ranks choose the lottery &  & Appendix: "the higher-rank workers would prefer the downgrade while the low-rank workers may prefer the lottery" & No functional form of $u$, no equilibrium concept and no transition kernel are defined. M13 gives the risk-neutral version only & --- & \texttt{OPEN} \\ \midrule
M21 & Mobility falls in $b$. Lottery play rises in $a$ &  & Appendix: "higher costs (b) discourage mobility and higher stakes (a) encourage bets" & Mobility measure undefined. The second half presumes higher $a$ means higher stakes, which M5 refutes for the printed ladder & --- & \texttt{OPEN} \\ \midrule
M22 & Upward mobility falls as a worker upskills &  & Appendix: "The more a worker upskills oneself, the harder it becomes for the worker to move up (as the swap pool gets limited)" & No formal object links upskilling to the swap pool & --- & \texttt{OPEN} \\ \midrule
M23 & A long-run transition kernel on ranks follows from the cost function &  & Appendix: "The long-term transition-probability can thus be determined from the costs of moving up the ladder" & No kernel is defined & --- & \texttt{OPEN} \\ \midrule
\end{longtable}}

\section{Literature claims}
{\small
\begin{longtable}{>{\raggedright\arraybackslash}p{0.8cm} >{\raggedright\arraybackslash}p{4.6cm} >{\raggedright\arraybackslash}p{6.6cm} >{\raggedright\arraybackslash}p{5.2cm} >{\raggedright\arraybackslash}p{3.4cm} >{\raggedright\arraybackslash}p{2.5cm}}
\toprule ID & Source & Claim & Evidence seen & Verbatim quote & Status \\ \midrule \endhead
L1 & Cunha, F. and Heckman, J. J. (2007). The Technology of Skill Formation. American Economic Review 97(2), 31--47 & Skill formation is written $\theta_{t+1}=f_t(h,\theta_t,I_t)$, with $\theta$ the skill stock, $h$ parental characteristics and $I$ investment. Self-productivity is $\partial f_t/\partial\theta_t>0$. Dynamic complementarity is $\partial^2 f_t/\partial\theta_t\partial I_t>0$ & NBER WP 12840 full text via web tool. Citation from IZA DP 2550 imprint\newline \url{https://www.nber.org/papers/w12840} & --- & \texttt{PRIMARY\_WEB} \\ \midrule
L2 & Gibbons, R. and Waldman, M. (1999). A Theory of Wage and Promotion Dynamics Inside Firms. Quarterly Journal of Economics 114(4), 1321--1358 & The model combines job assignment, on-the-job human capital acquisition and learning & NBER WP 6454 abstract page\newline \url{https://www.nber.org/papers/w6454} & --- & \texttt{ABSTRACT} \\ \midrule
L3 & Loury, G. C. (1977). A Dynamic Theory of Racial Income Differences. In Wallace, P. A. and LaMond, A. M. (eds.), Women, Minorities, and Employment Discrimination, Lexington Books, 153--186 & Under racial social grouping, children with equal family income face different community environments and different costs of becoming skilled. Equal opportunity alone need not remove an initial racial earnings gap & University of Chicago teaching handout summarising the chapter & --- & \texttt{SECONDARY} \\ \midrule
L4 & Borjas, G. J. (1992). Ethnic Capital and Intergenerational Mobility. Quarterly Journal of Economics 107(1), 123--150 & A child's skills depend on parental skills and on the average skill of the parents' ethnic group in the parents' generation & NBER WP 3788 abstract page\newline \url{https://gborjas.scholars.harvard.edu/resource/qje1992pdf} & --- & \texttt{ABSTRACT} \\ \midrule
L5 & B\'enabou, R. (1996). Equity and Efficiency in Human Capital Investment: The Local Connection. Review of Economic Studies 63(2), 237--264 & Small differences in endowments, preferences or credit access can produce large polarisation through community formation. Equilibrium stratification tends to be excessive & IDEAS abstract page. Full text paywalled & --- & \texttt{ABSTRACT} \\ \midrule
L6 & Bowles, S., Loury, G. C. and Sethi, R. (2014). Group Inequality. Journal of the European Economic Association 12(1), 129--152 & Skill cost $c(a,\sigma)$ falls in ability and in peer skill $\sigma_i=\eta s_i+(1-\eta)s$, where $\eta$ is the own-group share of social affiliates. Group inequality can persist under identical ability distributions only if segregation $\eta$ is high enough. Stability of the equal state changes at a threshold $\hat\eta$ (Proposition 3) & Full text on G. Loury's Brown University page via web tool\newline \url{https://www.brown.edu/Departments/Economics/Faculty/Glenn_Loury/louryhomepage/cvandbio/GROUP%20INEQUALITY%20-%20Bowles%20-%202014%20-%20Journal%20of%20the%20European%20Economic%20Association%20-%20Wiley%20Online%20Library.pdf} & --- & \texttt{PRIMARY\_WEB} \\ \midrule
L7 & Calv\'o-Armengol, A. and Jackson, M. O. (2004). The Effects of Social Networks on Employment and Inequality. American Economic Review 94(3), 426--453 (IDEAS lists 426--454) & Job information arrives at rate $a$ and is passed by employed agents to unemployed contacts. Long-run employment of path-connected agents is positively correlated (Proposition 1). Re-employment probability falls with spell length (Proposition 3). Groups with identical networks but worse initial states have persistently lower employment (Proposition 4) & Full text on a course page (NYU) via web tool\newline \url{https://pages.nyu.edu/debraj/Courses/NewRes05/Papers/CalvoJackson.pdf} & --- & \texttt{PRIMARY\_WEB} \\ \midrule
L8 & Rosen, S. (1986). Prizes and Incentives in Elimination Tournaments. American Economic Review 76(4), 701--715 & In sequential elimination tournaments the top prize must be raised to sustain effort in later rounds & IDEAS abstract page\newline \url{https://www.nber.org/papers/w1668} & --- & \texttt{ABSTRACT} \\ \midrule
L9 & Ghiglino, C. and Goyal, S. (2010). Keeping Up with the Neighbors: Social Interaction in a Market Economy. Journal of the European Economic Association 8(1), 90--119 & Utility falls in neighbours' consumption. Equilibrium consumption depends on network centrality. Integrating segregated communities lowers the utility of the poorer & IDEAS abstract page. Full text paywalled & --- & \texttt{ABSTRACT} \\ \midrule
L10 & Immorlica, N., Kranton, R., Manea, M. and Stoddard, G. (2017). Social Status in Networks. American Economic Journal: Microeconomics 9(1), 1--30 & Equilibria partition agents into classes. Network cohesion determines each class's action. The highest class is the largest maximally cohesive set & Duke Scholars abstract page. Definition of status not seen & --- & \texttt{ABSTRACT} \\ \midrule
L11 & Luttmer, E. F. P. (2005). Neighbors as Negatives: Relative Earnings and Well-Being. Quarterly Journal of Economics 120(3), 963--1002 & Holding own income fixed, higher local average earnings are associated with lower self-reported happiness & NBER WP 10667 abstract page\newline \url{https://www.nber.org/papers/w10667} & --- & \texttt{ABSTRACT} \\ \midrule
L12 & Altonji, J. G. and Pierret, C. R. (2001). Employer Learning and Statistical Discrimination. Quarterly Journal of Economics 116(1), 313--350 & The wage return to hard-to-observe productivity rises with experience and the return to education falls. Firms do not fully use the information in race & NBER WP 6279 abstract page\newline \url{https://www.nber.org/papers/w6279} & --- & \texttt{ABSTRACT} \\ \midrule
L13 & Fershtman, C., Murphy, K. M. and Weiss, Y. (1996). Social Status, Education, and Growth. Journal of Political Economy 104(1), 108--132 & When status matters, low-ability high-wealth individuals can crowd high-ability low-wealth individuals out of schooling, which lowers growth. A more equal wealth distribution can raise growth by reducing demand for status & IDEAS abstract page of the Tel Aviv working paper & --- & \texttt{ABSTRACT} \\ \midrule
\end{longtable}}

\section{Concluding remarks}
No concluding claims are admitted in this pass.

\section{References}
{\small
\begin{longtable}{>{\raggedright\arraybackslash}p{2.7cm} >{\raggedright\arraybackslash}p{5.0cm} >{\raggedright\arraybackslash}p{5.4cm} >{\raggedright\arraybackslash}p{5.4cm} >{\raggedright\arraybackslash}p{1.8cm} >{\raggedright\arraybackslash}p{2.6cm}}
\toprule Key & Full citation & Why the manuscript depends on it & Impact if not cited & Cited by & Maths / Lean \\ \midrule \endhead
\nolinkurl{CunhaHeckman2007} & Cunha, F. and Heckman, J. J. (2007). The Technology of Skill Formation. American Economic Review 97(2), 31--47 & Candidate law of motion for the human capital stock $h$ in the micro state (TODO: state space) & No current claim rests on this paper. Without it the law of motion for $h$ has no cited form and must be posited & none & yes / \texttt{NOT\_STARTED} \\ \midrule
\nolinkurl{GibbonsWaldman1999} & Gibbons, R. and Waldman, M. (1999). A Theory of Wage and Promotion Dynamics Inside Firms. Quarterly Journal of Economics 114(4), 1321--1358 & Precedent for a state that combines a job ladder with on-the-job human capital accumulation (TODO: state space) & No current claim rests on this paper. Without it the joint rank and human capital state has no cited precedent & none & yes / \texttt{NOT\_STARTED} \\ \midrule
\nolinkurl{Loury1977} & Loury, G. C. (1977). A Dynamic Theory of Racial Income Differences. In Wallace, P. A. and LaMond, A. M. (eds.), Women, Minorities, and Employment Discrimination, Lexington Books, 153--186 & Precedent for a group-level effect on the cost of acquiring skill (TODO: where categorical friction enters) & No current claim rests on this paper. Without it the route by which categorical friction enters skill acquisition has no cited origin & none & to confirm / \texttt{NOT\_STARTED} \\ \midrule
\nolinkurl{Borjas1992} & Borjas, G. J. (1992). Ethnic Capital and Intergenerational Mobility. Quarterly Journal of Economics 107(1), 123--150 & Precedent for group-average skill entering an individual's skill formation (TODO: where categorical friction enters) & No current claim rests on this paper. Without it group capital in the law of motion for $h$ has no cited basis & none & to confirm / \texttt{NOT\_STARTED} \\ \midrule
\nolinkurl{Benabou1996} & B\'enabou, R. (1996). Equity and Efficiency in Human Capital Investment: The Local Connection. Review of Economic Studies 63(2), 237--264 & Precedent for local spillovers and stratification amplifying small differences (TODO: locality) & No current claim rests on this paper. Without it the locality assumption has one fewer cited precedent & none & yes / \texttt{NOT\_STARTED} \\ \midrule
\nolinkurl{BowlesLourySethi2014} & Bowles, S., Loury, G. C. and Sethi, R. (2014). Group Inequality. Journal of the European Economic Association 12(1), 129--152 & Candidate operationalisation of categorical friction as an own-group affiliation share in peer skill (TODO: where categorical friction enters) & No current claim rests on this paper. Without it categorical friction has no cited functional form & none & yes / \texttt{NOT\_STARTED} \\ \midrule
\nolinkurl{CalvoArmengolJackson2004} & Calv\'o-Armengol, A. and Jackson, M. O. (2004). The Effects of Social Networks on Employment and Inequality. American Economic Review 94(3), 426--453 (IDEAS lists 426--454) & Candidate propagation channel for the rank kernel: opportunities passed through contacts (TODO: rank kernel) & No current claim rests on this paper. Without it the propagation channel of the kernel has no cited mechanism & none & yes / \texttt{NOT\_STARTED} \\ \midrule
\nolinkurl{Rosen1986} & Rosen, S. (1986). Prizes and Incentives in Elimination Tournaments. American Economic Review 76(4), 701--715 & Candidate rank-local kernel through sequential pairwise contests (TODO: rank kernel) & No current claim rests on this paper. Without it rank-local transitions have no cited contest structure & none & yes / \texttt{NOT\_STARTED} \\ \midrule
\nolinkurl{GhiglinoGoyal2010} & Ghiglino, C. and Goyal, S. (2010). Keeping Up with the Neighbors: Social Interaction in a Market Economy. Journal of the European Economic Association 8(1), 90--119 & Candidate reference-set rank: payoffs depend on neighbours rather than on the whole population (TODO: rank kernel) & No current claim rests on this paper. Without it reference-set rank has one fewer cited precedent & none & yes / \texttt{NOT\_STARTED} \\ \midrule
\nolinkurl{ImmorlicaEtAl2017} & Immorlica, N., Kranton, R., Manea, M. and Stoddard, G. (2017). Social Status in Networks. American Economic Journal: Microeconomics 9(1), 1--30 & Candidate reference-set rank through status in networks (TODO: rank kernel) & No current claim rests on this paper. Without it reference-set rank has one fewer cited precedent & none & yes / \texttt{NOT\_STARTED} \\ \midrule
\nolinkurl{Luttmer2005} & Luttmer, E. F. P. (2005). Neighbors as Negatives: Relative Earnings and Well-Being. Quarterly Journal of Economics 120(3), 963--1002 & Empirical motivation that comparisons are local (TODO: locality) & No current claim rests on this paper. Without it the locality assumption has no cited empirical motivation & none & no / \texttt{NOT\_APPLICABLE} \\ \midrule
\nolinkurl{AltonjiPierret2001} & Altonji, J. G. and Pierret, C. R. (2001). Employer Learning and Statistical Discrimination. Quarterly Journal of Economics 116(1), 313--350 & Candidate observer-learning channel, and evidence that employers do not fully use race information (TODO: rank kernel, categorical friction) & No current claim rests on this paper. Without it the observer-learning channel and the premise that firms do not act on categorisation have no cited basis & none & to confirm / \texttt{NOT\_STARTED} \\ \midrule
\nolinkurl{FershtmanMurphyWeiss1996} & Fershtman, C., Murphy, K. M. and Weiss, Y. (1996). Social Status, Education, and Growth. Journal of Political Economy 104(1), 108--132 & Positioning of the inequality and growth claim against a status-and-schooling model (deferred with the macro model) & No current claim rests on this paper. Without it the growth positioning has one fewer cited comparator & none & yes / \texttt{NOT\_STARTED} \\ \midrule
\end{longtable}}

\appendix
\section{Lean source}
\subsection*{Listing A.1: \nolinkurl{lean/IscLean.lean}}
\VerbatimInput[numbers=left,numbersep=4pt,fontsize=\footnotesize,frame=single]{lean/IscLean.lean}

\subsection*{Listing A.2: \nolinkurl{lean/IscLean/Lottery.lean}}
\VerbatimInput[numbers=left,numbersep=4pt,fontsize=\footnotesize,frame=single]{lean/IscLean/Lottery.lean}

\subsection*{Listing A.3: \nolinkurl{lean/IscLean/Ladder.lean}}
\VerbatimInput[numbers=left,numbersep=4pt,fontsize=\footnotesize,frame=single]{lean/IscLean/Ladder.lean}

\subsection*{Listing A.4: \nolinkurl{lean/IscLean/Gap.lean}}
\VerbatimInput[numbers=left,numbersep=4pt,fontsize=\footnotesize,frame=single]{lean/IscLean/Gap.lean}

\section{Lean build files}
\subsection*{Listing B.1: \nolinkurl{lean/lean-toolchain}}
\VerbatimInput[numbers=left,numbersep=4pt,fontsize=\scriptsize,frame=single]{lean/lean-toolchain}

\subsection*{Listing B.2: \nolinkurl{lean/lakefile.toml}}
\VerbatimInput[numbers=left,numbersep=4pt,fontsize=\scriptsize,frame=single]{lean/lakefile.toml}

\subsection*{Listing B.3: \nolinkurl{lean/lake-manifest.json}}
\VerbatimInput[numbers=left,numbersep=4pt,fontsize=\scriptsize,frame=single]{lean/lake-manifest.json}

\section{SymPy evidence}
\subsection*{Listing C.1: \nolinkurl{sympy/micro_checks.py}}
\VerbatimInput[numbers=left,numbersep=4pt,fontsize=\footnotesize,frame=single]{sympy/micro_checks.py}

\section{Measurement-map rows cited by headline and concluding claims}
No headline or concluding claim cites a measurement-map row in this pass.
\end{document}
